\section{Reconstrucción reproducible en Python}

\subsection{Propósito de la reconstrucción}

La reconstrucción en Python constituye la siguiente etapa de la genealogía computacional de la planta PMSM. Después de establecer la máquina de referencia, su parametrización, la formulación matemática y la implementación MATLAB/Simscape, se desarrolla una representación explícita, versionable y reproducible mediante código Python.

El objetivo no es reproducir literalmente la estructura interna de MATLAB/Simscape. La finalidad es expresar de manera explícita las relaciones físicas que pueden justificarse mediante la documentación, la formulación matemática y la evidencia recuperada del proyecto.

La reconstrucción puede representarse como:

\begin{equation}
\text{referencia MATLAB/Simscape}
\rightarrow
\text{formulación física}
\rightarrow
\text{Python explícito}.
\end{equation}

Esta estrategia permite separar la evidencia histórica de la implementación actualmente mantenida.

La implementación Python constituye además la semilla computacional sobre la cual se desarrollarán posteriormente los mapas electromagnéticos, los modelos de mayor fidelidad y las diferentes estrategias de control que utilizarán la planta como plataforma común.

\subsection{Motivación de la migración a Python}

La utilización de Python proporciona características especialmente útiles para una reconstrucción reproducible:

\begin{itemize}
    \item código fuente versionable;
    \item parametrización explícita;
    \item experimentos automatizables;
    \item generación reproducible de datos;
    \item integración con herramientas de análisis numérico;
    \item construcción sistemática de mapas;
    \item generación automática de figuras y tablas;
    \item trazabilidad entre código, parámetros y resultados.
\end{itemize}

La migración permite, además, desacoplar la definición de la planta de una determinada herramienta gráfica.

La planta puede ejecutarse desde scripts y módulos Python, facilitando que diferentes experimentos utilicen exactamente la misma implementación base.

La decisión de utilizar Python no implica eliminar MATLAB/Simscape. Ambas plataformas cumplen funciones diferentes dentro de la arquitectura del proyecto.

\subsection{Relación entre referencia y reconstrucción}

La relación entre ambas implementaciones debe entenderse en términos de evidencia y formulación, no de copia estructural.

La referencia MATLAB/Simscape aporta información sobre:

\begin{itemize}
    \item representación física;
    \item conexiones entre dominios;
    \item condiciones de operación;
    \item variables observables;
    \item arquitectura histórica;
    \item resultados de simulación.
\end{itemize}

La reconstrucción Python convierte la información físicamente justificable en una implementación explícita.

Por tanto:

\begin{equation}
\text{Simscape}
\neq
\text{código fuente de Python}.
\end{equation}

La relación correcta es:

\begin{equation}
\text{evidencia}
+
\text{formulación}
+
\text{parametrización}
\rightarrow
\text{implementación Python}.
\end{equation}

Esta distinción evita atribuir a Python comportamientos que todavía no hayan sido demostrados o trazados.

\subsection{Arquitectura del código}

La implementación activa se organiza alrededor del paquete:

\begin{equation}
\texttt{pmsm\_hybrid\_control}.
\end{equation}

Dentro de esta estructura, la planta se separa de las capas de control.

La cadena principal de la planta es:

\begin{equation}
\texttt{parameters.py}
\rightarrow
\texttt{PMSMParameters}
\rightarrow
\texttt{PMSMBase}.
\end{equation}

Las interfaces de entrada y salida se representan mediante:

\begin{equation}
\texttt{PlantInput}
\rightarrow
\texttt{PMSMBase}
\rightarrow
\texttt{PlantOutput}.
\end{equation}

Esta separación permite que la planta pueda ser utilizada por diferentes experimentos sin incorporar directamente la lógica del controlador.

\subsection{Fuente de parámetros}

El módulo \texttt{parameters.py} establece la fuente de verdad para la parametrización M4 utilizada por la planta actual.

El propio módulo documenta que la fuente principal es:

\begin{equation}
\texttt{01\_M4\_HARD\_DATA.md}.
\end{equation}

Los parámetros se acompañan de información de trazabilidad.

La estructura utilizada permite asociar a un valor:

\begin{equation}
\{
\text{valor},
\text{estado},
\text{fuente},
\text{unidad},
\text{nota}
\}.
\end{equation}

Esta representación evita que un número aislado pierda su procedencia al ser utilizado en diferentes partes del proyecto.

\subsection{Estados de trazabilidad}

La implementación utiliza estados explícitos para distinguir el nivel de evidencia de los parámetros.

Entre ellos se encuentran:

\begin{itemize}
    \item \texttt{CONFIRMED};
    \item \texttt{PROPOSED};
    \item \texttt{PENDING\_VERIFICATION}.
\end{itemize}

Estos estados no representan niveles de calidad numérica, sino estados de conocimiento y procedencia.

Un valor confirmado puede utilizarse como parte de la parametrización establecida.

Un valor propuesto representa una decisión de modelado que todavía puede evolucionar.

Un valor pendiente no debe incorporarse silenciosamente como si fuera un parámetro físicamente confirmado.

\subsection{Parámetros M4 de la planta}

La parametrización concentrada actualmente utilizada contiene:

\begin{equation}
R_s=0.05\,\Omega,
\end{equation}

\begin{equation}
L_d=1.2\times10^{-3}\,\mathrm{H},
\end{equation}

\begin{equation}
L_q=1.2\times10^{-3}\,\mathrm{H},
\end{equation}

\begin{equation}
\lambda_m=0.107245\,\mathrm{Wb},
\end{equation}

\begin{equation}
p=4,
\end{equation}

y:

\begin{equation}
J=0.00416\,\mathrm{kg\,m^2}.
\end{equation}

Estos valores constituyen la parametrización M4 actualmente adoptada.

La igualdad:

\begin{equation}
L_d=L_q
\end{equation}

corresponde a la representación concentrada utilizada actualmente. No debe interpretarse como una afirmación universal de que la máquina real carezca de cualquier anisotropía.

\subsection{Flujo magnético permanente}

El valor utilizado actualmente para el enlace de flujo permanente es:

\begin{equation}
\lambda_m=0.107245\,\mathrm{Wb}.
\end{equation}

Este valor se obtuvo a partir de la constante de par documentada y del número de pares de polos:

\begin{equation}
\lambda_m=
\frac{k_T}
{(3/2)\sqrt{2}\,p}.
\end{equation}

Utilizando:

\begin{equation}
k_T=0.91\,\mathrm{N\,m/A}
\end{equation}

y:

\begin{equation}
p=4,
\end{equation}

se obtiene aproximadamente el valor utilizado por M4.

La procedencia numérica y la procedencia documental deben mantenerse diferenciadas. El valor se encuentra incorporado a la parametrización M4, mientras que su trazabilidad documental detallada permanece identificada para futura verificación.

\subsection{Separación de parámetros históricos}

Durante la reconstrucción se encontraron valores pertenecientes a implementaciones anteriores.

Un caso relevante es:

\begin{equation}
\lambda_m=0.1100\,\mathrm{Wb}.
\end{equation}

Este valor aparece en configuraciones históricas del FOC reconstruido.

No se utiliza como sustituto del valor M4 de la planta.

La separación se expresa como:

\begin{equation}
\lambda_{m,\mathrm{plant}}
\neq
\lambda_{m,\mathrm{FOC,historical}}.
\end{equation}

La diferencia no se corrige silenciosamente porque forma parte de la trazabilidad del proceso de recuperación.

La misma regla se aplica a otros parámetros históricos.

\subsection{Parámetros mecánicos}

La inercia actualmente utilizada es:

\begin{equation}
J=0.00416\,\mathrm{kg\,m^2}.
\end{equation}

Este valor corresponde a la configuración con freno considerada en la parametrización M4.

La documentación del fabricante también proporciona una inercia de aproximadamente:

\begin{equation}
J=0.00274\,\mathrm{kg\,m^2}
\end{equation}

para la configuración sin freno.

La selección del valor M4 debe permanecer documentada porque la inercia depende de la configuración mecánica considerada.

\subsection{Parámetros todavía pendientes}

La reconstrucción no incorpora como parámetros activos de M4 determinados valores históricos de fricción.

En particular:

\begin{equation}
B
\end{equation}

y:

\begin{equation}
T_{\mathrm{fric}}
\end{equation}

permanecen pendientes de verificación.

Valores históricos asociados a CloudE no se trasladan automáticamente a la planta actual.

Esta decisión permite mantener una separación clara entre:

\begin{equation}
\text{parámetro confirmado}
\end{equation}

y:

\begin{equation}
\text{valor histórico recuperado}.
\end{equation}

\subsection{Clase \texttt{PMSMParameters}}

La clase \texttt{PMSMParameters} concentra los parámetros necesarios para ejecutar la planta M4.

Su función es evitar que las ecuaciones dependan de valores dispersos en diferentes módulos.

Conceptualmente:

\begin{equation}
\Theta=
\{
R_s,
L_d,
L_q,
\lambda_m,
p,
J
\}.
\end{equation}

El objeto de parámetros representa este conjunto como una unidad.

La planta recibe así una configuración explícita y verificable.

Esta estructura facilita la creación de diferentes configuraciones futuras sin modificar las ecuaciones fundamentales.

\subsection{Validación de parámetros}

La implementación incorpora validaciones básicas antes de ejecutar la planta.

Entre las condiciones comprobadas se encuentran:

\begin{itemize}
    \item resistencia positiva;
    \item inductancias positivas;
    \item flujo permanente positivo;
    \item inercia positiva;
    \item número de pares de polos consistente;
    \item valores numéricos válidos.
\end{itemize}

Estas comprobaciones no constituyen una validación física completa de la máquina.

Su finalidad es impedir que errores elementales de parametrización pasen inadvertidos durante la ejecución.

\subsection{Estados de la planta}

La clase \texttt{PMSMBase} utiliza cuatro estados principales:

\begin{equation}
x=
\begin{bmatrix}
i_d&
i_q&
\omega_m&
\theta_m
\end{bmatrix}^{T}.
\end{equation}

Los dos primeros representan las corrientes en el marco $dq$.

El tercero representa la velocidad mecánica del rotor.

El cuarto representa la posición mecánica.

Estos estados corresponden directamente a la formulación matemática desarrollada anteriormente.

\subsection{Entradas de la planta}

La interfaz \texttt{PlantInput} representa las magnitudes que la implementación actual utiliza para avanzar la planta.

Incluye:

\begin{equation}
v_a,\quad v_b,\quad v_c,
\end{equation}

el par de carga:

\begin{equation}
T_{\mathrm{load}},
\end{equation}

y la posición mecánica utilizada por la implementación:

\begin{equation}
\theta_m.
\end{equation}

La presencia de $\theta_m$ como entrada debe entenderse como una característica de la implementación actual y no como una afirmación de que la posición del rotor sea una entrada física independiente de la máquina.

La consolidación futura de esta interfaz deberá distinguir de forma más estricta entre estado interno y señal externa de posición.

\subsection{Salidas de la planta}

La interfaz \texttt{PlantOutput} expone:

\begin{equation}
i_a,\quad i_b,\quad i_c,
\end{equation}

\begin{equation}
i_d,\quad i_q,
\end{equation}

\begin{equation}
\omega_m,\quad \theta_m,
\end{equation}

y:

\begin{equation}
T_e.
\end{equation}

Esta interfaz permite que la misma planta sea utilizada por diferentes capas de control y análisis.

La salida no contiene directamente parámetros del controlador.

La separación permite utilizar la planta como componente independiente.

\subsection{Transformación de variables trifásicas}

La implementación recibe variables trifásicas y las transforma al marco de referencia utilizado por las ecuaciones.

La secuencia conceptual es:

\begin{equation}
abc
\rightarrow
\alpha\beta
\rightarrow
dq.
\end{equation}

La implementación actual utiliza una transformación de Clarke basada en:

\begin{equation}
i_\alpha=i_a,
\end{equation}

y:

\begin{equation}
i_\beta=
\frac{i_b-i_c}{\sqrt{3}}.
\end{equation}

Posteriormente se aplica la transformación de Park utilizando la posición eléctrica.

La transformación inversa permite reconstruir las variables trifásicas de salida.

\subsection{Posición eléctrica}

La posición eléctrica se obtiene mediante:

\begin{equation}
\theta_e=p\theta_m.
\end{equation}

De forma equivalente, la velocidad eléctrica es:

\begin{equation}
\omega_e=p\omega_m.
\end{equation}

Esta relación constituye uno de los principales acoplamientos entre los dominios mecánico y eléctrico.

El número utilizado actualmente es:

\begin{equation}
p=4,
\end{equation}

correspondiente a ocho polos físicos.

\subsection{Ecuaciones eléctricas implementadas}

La dinámica eléctrica utilizada por \texttt{PMSMBase} es:

\begin{equation}
\frac{di_d}{dt}
=
\frac{
v_d-R_si_d+\omega_eL_qi_q
}
{L_d},
\end{equation}

y:

\begin{equation}
\frac{di_q}{dt}
=
\frac{
v_q-R_si_q-\omega_e(L_di_d+\lambda_m)
}
{L_q}.
\end{equation}

Estas ecuaciones corresponden directamente a la formulación matemática establecida para el modelo concentrado.

La implementación no introduce en esta etapa dependencias angulares de inductancia ni mapas electromagnéticos.

\subsection{Par electromagnético}

El par electromagnético utilizado como relación de referencia es:

\begin{equation}
T_e=
\frac{3}{2}p
\left[
\lambda_mi_q+
(L_d-L_q)i_di_q
\right].
\end{equation}

Dado que el modelo M4 utiliza:

\begin{equation}
L_d=L_q,
\end{equation}

la expresión se reduce a:

\begin{equation}
T_e=
\frac{3}{2}p\lambda_mi_q.
\end{equation}

Esta ecuación constituye la relación de par utilizada por la implementación M4. No se presenta como una reproducción literal de una ley interna de Simscape.

\subsection{Dinámica mecánica}

La velocidad mecánica se integra mediante:

\begin{equation}
\frac{d\omega_m}{dt}
=
\frac{T_e-T_{\mathrm{load}}}{J}.
\end{equation}

La posición se obtiene mediante:

\begin{equation}
\frac{d\theta_m}{dt}
=
\omega_m.
\end{equation}

El modelo actual no introduce explícitamente términos separados de fricción viscosa o par de fricción porque dichos parámetros permanecen pendientes de verificación.

Por tanto, $T_{\mathrm{load}}$ representa el par resistente externo utilizado en el experimento.

\subsection{Integración numérica}

La implementación actual utiliza integración explícita de Euler.

Para un estado genérico:

\begin{equation}
x_{k+1}
=
x_k+\Delta t\,\dot{x}_k.
\end{equation}

Aplicado a las corrientes:

\begin{equation}
i_{d,k+1}
=
i_{d,k}
+
\Delta t
\frac{di_d}{dt},
\end{equation}

\begin{equation}
i_{q,k+1}
=
i_{q,k}
+
\Delta t
\frac{di_q}{dt}.
\end{equation}

Para la velocidad:

\begin{equation}
\omega_{m,k+1}
=
\omega_{m,k}
+
\Delta t
\frac{T_e-T_{\mathrm{load}}}{J}.
\end{equation}

La posición se actualiza mediante:

\begin{equation}
\theta_{m,k+1}
=
\theta_{m,k}
+
\Delta t\,\omega_{m,k}.
\end{equation}

El método se mantiene explícito porque facilita la inspección del comportamiento numérico y la reproducibilidad del experimento.

\subsection{Paso temporal de referencia}

El experimento base utiliza:

\begin{equation}
\Delta t=50\,\mu\mathrm{s}.
\end{equation}

Esto corresponde a una frecuencia de integración de:

\begin{equation}
f_s=20\,\mathrm{kHz}.
\end{equation}

El paso temporal constituye una condición numérica del experimento y no un parámetro físico de la máquina.

Esta distinción permite cambiar el integrador o realizar estudios de convergencia sin alterar la parametrización física.

\subsection{Condiciones iniciales}

La implementación base inicializa las variables de estado en cero.

Conceptualmente:

\begin{equation}
i_d(0)=0,
\end{equation}

\begin{equation}
i_q(0)=0,
\end{equation}

\begin{equation}
\omega_m(0)=0,
\end{equation}

y:

\begin{equation}
\theta_m(0)=0.
\end{equation}

Estas condiciones corresponden al experimento de arranque utilizado como referencia.

Una condición inicial diferente constituye una configuración experimental diferente y debe quedar registrada junto con los resultados.

\subsection{Peculiaridad actual de \texorpdfstring{$\theta_m$}{theta m}}

La implementación actual presenta una característica que debe conservarse explícitamente en la trazabilidad.

\texttt{PlantInput} contiene $\theta_m$, mientras que \texttt{PMSMBase} también mantiene una variable interna para la posición.

Durante el paso de simulación, la posición recibida se utiliza para calcular la posición eléctrica y se asigna al estado interno antes de la actualización dinámica.

Esto implica que actualmente existe una mezcla entre:

\begin{equation}
\text{posición externa}
\end{equation}

y:

\begin{equation}
\text{posición integrada internamente}.
\end{equation}

La condición no se modifica silenciosamente porque forma parte del estado actual del código.

Debe considerarse una tarea de consolidación futura de la interfaz de la planta.

\subsection{Ausencia de un solver FEM}

La implementación \texttt{PMSMBase} no constituye un solver de elementos finitos.

Su nivel de fidelidad corresponde a un modelo concentrado basado en parámetros.

La estructura es:

\begin{equation}
\text{PMSMBase}
=
\text{modelo matemático parametrizado}.
\end{equation}

No debe afirmarse que la implementación actual reproduce por sí misma la distribución espacial del campo magnético.

La dependencia angular detallada, las saturaciones locales y los armónicos espaciales requieren una representación de mayor fidelidad.

\subsection{Ausencia de mapas electromagnéticos activos}

El repositorio contiene módulos destinados a una futura representación mediante mapas.

Sin embargo, la configuración M4 activa utiliza relaciones analíticas.

El flujo de referencia es:

\begin{equation}
\psi_d=L_di_d+\lambda_m,
\end{equation}

\begin{equation}
\psi_q=L_qi_q.
\end{equation}

Las tablas recuperadas de etapas anteriores permanecen separadas y marcadas como información recuperada pendiente de verificación semántica.

Por tanto, no se consideran automáticamente mapas validados de la máquina.

\subsection{Evolución hacia mapas}

Cuando se disponga de datos electromagnéticos suficientemente trazables, las relaciones analíticas podrán sustituirse o complementarse mediante funciones dependientes del punto de operación.

Una representación futura puede adoptar la forma:

\begin{equation}
\psi_d=
f(i_d,i_q,\theta_e),
\end{equation}

\begin{equation}
\psi_q=
g(i_d,i_q,\theta_e).
\end{equation}

De forma correspondiente, el par podrá evaluarse directamente desde el estado electromagnético.

La arquitectura actual permite introducir esta evolución sin redefinir la identidad de la planta.

\subsection{Evolución del nivel de fidelidad}

La reconstrucción Python se organiza alrededor de una evolución progresiva:

\begin{equation}
\mathrm{M4}
\rightarrow
\mathrm{mapas}
\rightarrow
\mathrm{modelo\ no\ lineal}
\rightarrow
\mathrm{FEM/ROM}
\rightarrow
\mathrm{Digital\ Twin}.
\end{equation}

Cada etapa debe incorporar nueva evidencia.

La etapa M4 proporciona una semilla reproducible.

Los mapas permitirán incorporar dependencias no representadas por el modelo concentrado.

El modelo no lineal podrá incorporar saturación y dependencia angular.

FEM y ROM podrán proporcionar información espacial o modelos reducidos derivados de representaciones de mayor fidelidad.

\subsection{Separación entre planta y controlador}

La reconstrucción Python mantiene separada la planta de las estrategias de control.

La planta responde a entradas eléctricas y mecánicas y produce variables físicas.

El FOC utiliza dichas variables para generar las referencias necesarias para el accionamiento.

La relación conceptual es:

\begin{equation}
\text{FOC}
\rightarrow
\text{accionamiento}
\rightarrow
\text{PMSMBase}.
\end{equation}

Y en realimentación:

\begin{equation}
\text{PMSMBase}
\rightarrow
\text{mediciones}
\rightarrow
\text{FOC}.
\end{equation}

Esta separación permite utilizar posteriormente la misma planta con otras estrategias.

\subsection{Separación de \texttt{foc} y \texttt{foc\_reconstructed}}

El repositorio contiene dos estructuras relacionadas con FOC.

El directorio:

\begin{equation}
\texttt{foc/}
\end{equation}

representa actualmente una estructura de interfaz o esqueleto pendiente de consolidación.

El directorio:

\begin{equation}
\texttt{foc\_reconstructed/}
\end{equation}

contiene el FOC funcional utilizado por los experimentos de integración actuales.

Esta distinción evita presentar una interfaz placeholder como si fuera el controlador funcional utilizado en los resultados.

\subsection{Parámetros del FOC reconstruido}

El FOC reconstruido utiliza parámetros que deben mantenerse separados de la planta.

Entre los valores históricos recuperados se encuentran:

\begin{equation}
K_{p,id}=K_{p,iq}=15.0796,
\end{equation}

\begin{equation}
K_{i,id}=K_{i,iq}=628.3185,
\end{equation}

\begin{equation}
K_{p,\omega}=5.2804,
\end{equation}

y:

\begin{equation}
K_{i,\omega}=1.2693.
\end{equation}

Estos valores describen una configuración particular del controlador.

No deben incorporarse a la clase de parámetros físicos de la PMSM.

\subsection{Desacoplamiento del FOC}

El FOC reconstruido incorpora términos de desacoplamiento basados en la formulación $dq$.

Los términos utilizados como referencia son:

\begin{equation}
v_{d,\mathrm{dec}}
=
-\omega_eL_qi_q,
\end{equation}

y:

\begin{equation}
v_{q,\mathrm{dec}}
=
\omega_e
(L_di_d+\lambda_m).
\end{equation}

Estos términos pertenecen al controlador.

Su existencia no modifica la ecuación física de la planta.

Esta separación es fundamental para que la misma planta pueda recibir diferentes estrategias de control.

\subsection{Límites del FOC}

El FOC reconstruido también incorpora límites de corriente y tensión.

El límite histórico de corriente $q$ utilizado en esta configuración es aproximadamente:

\begin{equation}
|i_q|\leq112.43\,\mathrm{A}.
\end{equation}

La tensión máxima se relaciona con el bus mediante:

\begin{equation}
V_{\max}
=
0.95
\frac{V_{DC}}{\sqrt{3}}.
\end{equation}

Estos límites pertenecen a la estrategia de accionamiento y control.

No deben interpretarse como sustitutos de los límites físicos documentados para la máquina.

\subsection{Arquitectura del experimento B0}

La primera referencia cerrada de la reconstrucción Python utiliza una cadena M3--M6--M4.

La arquitectura es:

\begin{equation}
\omega^\ast
\rightarrow
\mathrm{FOC}
\rightarrow
(v_d^\ast,v_q^\ast)
\rightarrow
\mathrm{M3}
\rightarrow
(v_a,v_b,v_c)
\rightarrow
\mathrm{M4}.
\end{equation}

La realimentación devuelve:

\begin{equation}
\mathrm{M4}
\rightarrow
(i_a,i_b,i_c,\omega_m,\theta_m)
\rightarrow
\mathrm{FOC}.
\end{equation}

Esta configuración constituye un experimento de referencia de la planta reconstruida.

\subsection{Condición de bus del experimento B0}

El experimento B0 utiliza:

\begin{equation}
V_{DC}=540\,\mathrm{V}.
\end{equation}

Este valor forma parte de la configuración experimental del accionamiento.

No se incorpora a \texttt{PMSMParameters} porque no constituye un parámetro intrínseco de la máquina.

La misma planta puede utilizarse posteriormente con otra condición de bus si se mantienen documentadas las nuevas condiciones.

\subsection{Referencia de velocidad de B0}

La referencia de velocidad del experimento B0 es:

\begin{equation}
\omega^\ast=157.079633\,\mathrm{rad/s}.
\end{equation}

La carga cambia de cero a:

\begin{equation}
T_{\mathrm{load}}=23\,\mathrm{N\,m}
\end{equation}

a los $0.05\,\mathrm{s}$.

El experimento tiene una duración total de:

\begin{equation}
t_{\mathrm{stop}}=20\,\mathrm{s}.
\end{equation}

Estas magnitudes constituyen la configuración del experimento y no deben confundirse con parámetros permanentes de la planta.

\subsection{Resultado de referencia B0}

La ejecución instrumentada del experimento B0 produjo, al final de la simulación, aproximadamente:

\begin{equation}
\omega_m=157.023729\,\mathrm{rad/s},
\end{equation}

con:

\begin{equation}
i_d=-1.122894\,\mathrm{A},
\end{equation}

\begin{equation}
i_q=35.743789\,\mathrm{A},
\end{equation}

y:

\begin{equation}
T_e=23.000056\,\mathrm{N\,m}.
\end{equation}

La referencia de corriente $q$ al final fue aproximadamente:

\begin{equation}
i_q^\ast=35.761422\,\mathrm{A}.
\end{equation}

Estos valores pertenecen a la configuración experimental B0 y no constituyen nuevos parámetros de la máquina.

\subsection{Métricas instrumentadas}

Entre las métricas obtenidas en B0 se encuentran:

\begin{equation}
e_{\omega,\mathrm{final}}
=
0.055903785\,\mathrm{rad/s},
\end{equation}

y un error relativo final aproximado de:

\begin{equation}
e_{\omega,\mathrm{rel}}
=
0.035589454\%.
\end{equation}

El tiempo de establecimiento al 1\% registrado fue aproximadamente:

\begin{equation}
t_s=6.12635\,\mathrm{s}.
\end{equation}

También se registraron:

\begin{equation}
|i_q|_{\max}=120.2163\,\mathrm{A},
\end{equation}

\begin{equation}
|i_d|_{\max}=3.39164\,\mathrm{A},
\end{equation}

y:

\begin{equation}
|V_{dq}|_{\max}=296.1843\,\mathrm{V}.
\end{equation}

La salida PWM permaneció dentro del intervalo:

\begin{equation}
0\leq D\leq1.
\end{equation}

Estas métricas constituyen resultados del experimento y no valores de parametrización física.

\subsection{Trazabilidad de los resultados B0}

La trayectoria experimental y las métricas fueron almacenadas con identificadores SHA-256.

La trayectoria B0 presenta el hash:

\begin{equation}
\texttt{68c89e3d88ed51d18dd9fe07061d1185169aa95c7dd6a43846d1ba173b914278}.
\end{equation}

El archivo de métricas presenta:

\begin{equation}
\texttt{c0f4fc2e56430b912089660173a18db8808bd6a42512e1691c977a68b1a002ea}.
\end{equation}

Estos identificadores permiten asociar un conjunto de resultados con una versión concreta de los datos generados.

El hash no demuestra por sí mismo la validez física del resultado. Su función es garantizar la identidad del artefacto almacenado.

\subsection{Reproducibilidad experimental}

La reproducibilidad requiere más que conservar un archivo de resultados.

Debe ser posible reconstruir la configuración utilizada mediante:

\begin{equation}
C=
\{
\Theta,
V_{DC},
\omega^\ast,
T_{\mathrm{load}},
x_0,
\Delta t,
t_{\mathrm{stop}},
\text{control}
\}.
\end{equation}

La ejecución de un experimento debe registrar estos elementos junto con sus resultados.

De esta forma, un resultado puede ser asociado a:

\begin{equation}
\text{código}
+
\text{parámetros}
+
\text{configuración}
+
\text{datos}.
\end{equation}

Esta asociación constituye una de las principales ventajas de la reconstrucción reproducible.

\subsection{Resultados y generación automática}

Los resultados experimentales no deben incorporarse manualmente al código fuente del reporte.

El flujo previsto es:

\begin{equation}
\text{experimento Python}
\rightarrow
\text{datos}
\rightarrow
\text{métricas}
\rightarrow
\text{generador}
\rightarrow
\text{LaTeX}.
\end{equation}

El reporte puede entonces incorporar valores y figuras generados automáticamente.

Esta metodología reduce el riesgo de copiar manualmente resultados y permite repetir el proceso cuando cambien las condiciones experimentales.

\subsection{Relación con el generador de datos LaTeX}

El directorio del reporte contiene un generador destinado a producir archivos consumibles por LaTeX.

El archivo:

\begin{equation}
\texttt{report/scripts/generate\_latex\_data.py}
\end{equation}

constituye la infraestructura para esta transferencia.

El archivo generado:

\begin{equation}
\texttt{report/datos/resultados\_pmsm.tex}
\end{equation}

debe contener macros derivadas de resultados experimentales.

La plantilla actual no debe interpretarse como una indicación de que todas las métricas del proyecto ya estén automatizadas.

La infraestructura está preparada para incorporar el experimento de referencia conforme se consoliden sus resultados.

\subsection{Reproducibilidad y control de versiones}

La implementación Python forma parte del repositorio Git del proyecto.

Esto permite asociar:

\begin{equation}
\text{código}
\leftrightarrow
\text{commit}
\leftrightarrow
\text{experimento}
\leftrightarrow
\text{resultado}.
\end{equation}

La identificación mediante Git complementa los hashes de los archivos experimentales.

Una ejecución reproducible debe conservar suficiente información para determinar qué versión del código produjo un resultado determinado.

\subsection{Separación entre implementación y evidencia histórica}

La reconstrucción no incorpora automáticamente todos los scripts encontrados durante la recuperación.

Los scripts diagnósticos antiguos, pruebas incompletas y artefactos históricos se consideran material de trazabilidad salvo que hayan sido seleccionados explícitamente como parte de la implementación activa.

Esta regla evita contaminar el modelo actual con decisiones pertenecientes a experimentos anteriores.

El mismo criterio se aplica al material ROM/FEM histórico.

\subsection{Relación con los artefactos recuperados}

El proceso de recuperación proporcionó diferentes fuentes de información.

Estas pueden clasificarse como:

\begin{itemize}
    \item documentación de fabricante;
    \item implementación MATLAB/Simscape;
    \item archivos históricos;
    \item parámetros M4;
    \item código Python reconstruido;
    \item resultados experimentales.
\end{itemize}

Cada fuente tiene una función diferente.

La reconstrucción Python no intenta resolver automáticamente las contradicciones entre fuentes.

Cuando existen discrepancias, estas deben conservarse y clasificarse hasta que exista evidencia suficiente para resolverlas.

\subsection{Planta reproducible como semilla}

El modelo M4 constituye una semilla computacional porque permite ejecutar experimentos sin depender de un modelo gráfico externo.

La semilla contiene:

\begin{equation}
\{
\text{parámetros},
\text{ecuaciones},
\text{estados},
\text{entradas},
\text{salidas},
\text{integrador}
\}.
\end{equation}

A partir de esta base pueden construirse experimentos más complejos sin modificar necesariamente la definición fundamental de la planta.

Esto permite utilizar la misma planta como entorno común de evaluación.

\subsection{Preparación para control avanzado}

Una planta explícita y reproducible facilita el desarrollo de diferentes estrategias de control.

La arquitectura prevista es:

\begin{equation}
\text{PMSM}
\rightarrow
\left\{
\begin{array}{c}
\mathrm{FOC}\\
\mathrm{MPC}\\
\mathrm{ANFIS}\\
\mathrm{TinyML}
\end{array}
\right.
\end{equation}

La planta permanece común mientras las estrategias de control cambian.

Esta condición permite que las diferencias observadas entre estrategias puedan asociarse a sus respectivas decisiones de control, siempre que se mantengan constantes las condiciones de la planta y del experimento.

\subsection{FOC como controlador establecido}

En la arquitectura actual, FOC constituye el controlador funcional utilizado para cerrar el primer experimento de referencia.

Esto no significa que sea la única estrategia futura.

Su función actual es proporcionar una primera cadena cerrada que permita verificar la interacción:

\begin{equation}
\text{referencia}
\rightarrow
\text{control}
\rightarrow
\text{planta}
\rightarrow
\text{realimentación}.
\end{equation}

Las estrategias MPC, ANFIS y TinyML se desarrollarán posteriormente sobre la misma base de planta y respetando las interfaces establecidas.

\subsection{Relación con MPC}

MPC requiere una representación de la planta que pueda utilizarse para predecir su evolución bajo diferentes acciones de control.

La reconstrucción Python proporciona una base para esa tarea.

Sin embargo, la existencia del modelo M4 no implica que el MPC ya esté implementado o validado.

La relación actual es preparatoria:

\begin{equation}
\text{M4 reproducible}
\rightarrow
\text{modelo predictivo}
\rightarrow
\text{MPC}.
\end{equation}

El modelo predictivo deberá conservar la trazabilidad de las simplificaciones adoptadas.

\subsection{Relación con ANFIS y neurodifuso}

La arquitectura prevista coloca ANFIS como capa supervisora o adaptativa y no como sustituto de la planta ni del controlador principal.

La separación conceptual es:

\begin{equation}
\text{ANFIS}
\rightarrow
\text{supervisión/adaptación}
\rightarrow
\text{FOC}.
\end{equation}

La planta permanece:

\begin{equation}
\text{PMSM}
\end{equation}

y el FOC continúa siendo el controlador funcional de la primera arquitectura cerrada.

La reconstrucción Python permite preparar los datos necesarios para futuras estrategias neurodifusas sin alterar la definición física de la planta.

\subsection{Relación con TinyML}

La misma filosofía se aplica a las futuras capas TinyML.

Los datos producidos por la planta pueden utilizarse para construir conjuntos de entrenamiento, validación y prueba.

La cadena futura puede representarse como:

\begin{equation}
\text{PMSM}
\rightarrow
\text{datos}
\rightarrow
\text{modelo TinyML}
\rightarrow
\text{supervisión}.
\end{equation}

La generación de estos datos debe conservar las condiciones físicas y experimentales utilizadas para producirlos.

De esta manera, la planta se convierte en una fuente controlada de escenarios para el desarrollo de inteligencia aplicada al control.

\subsection{Contratos de interfaz}

La reconstrucción Python mantiene la separación conceptual entre bloques mediante interfaces.

La planta no debe conocer la implementación interna del controlador.

El controlador tampoco debe depender de detalles internos de integración de la planta.

La comunicación conceptual es:

\begin{equation}
\text{controlador}
\rightarrow
[\text{contrato}]
\rightarrow
\text{planta},
\end{equation}

y:

\begin{equation}
\text{planta}
\rightarrow
[\text{contrato}]
\rightarrow
\text{controlador}.
\end{equation}

Esta separación será necesaria para trasladar posteriormente la arquitectura a Simulink/Simscape.

\subsection{Ventaja de la planta común}

Una planta común permite comparar diferentes estrategias sin reconstruir el sistema físico para cada controlador.

La condición puede expresarse como:

\begin{equation}
PMSM_{\mathrm{common}}
=
PMSM_{\mathrm{FOC}}
=
PMSM_{\mathrm{MPC}}
=
PMSM_{\mathrm{ANFIS}}
=
PMSM_{\mathrm{TinyML}}.
\end{equation}

La igualdad anterior representa la intención arquitectónica de utilizar una misma planta física de referencia, no la afirmación de que todos los modelos futuros tendrán exactamente el mismo nivel de fidelidad.

Conforme la planta evolucione, las diferentes estrategias deberán compararse sobre versiones identificadas de dicha planta.

\subsection{Control de fidelidad}

La evolución de la planta debe estar acompañada por una identificación clara del nivel de fidelidad.

El modelo M4 actual corresponde a una representación concentrada.

Las futuras versiones podrán incorporar:

\begin{itemize}
    \item mapas electromagnéticos;
    \item saturación;
    \item dependencia angular;
    \item armónicos;
    \item pérdidas adicionales;
    \item dependencia térmica;
    \item modelos FEM;
    \item modelos reducidos ROM.
\end{itemize}

Cada aumento de fidelidad debe conservar una relación explícita con la versión anterior.

Esto permite determinar si una mejora en una estrategia de control proviene realmente del controlador o de una modificación simultánea de la planta.

\subsection{Criterio de evolución}

La evolución del modelo no debe realizarse únicamente por complejidad.

Una representación de mayor fidelidad debe incorporar información adicional que pueda justificarse mediante datos o modelos de referencia.

El criterio conceptual es:

\begin{equation}
\text{mayor fidelidad}
\Rightarrow
\text{mayor información física verificable}.
\end{equation}

No:

\begin{equation}
\text{mayor complejidad}
\Rightarrow
\text{mayor validez}.
\end{equation}

Esta distinción es especialmente importante para evitar introducir parámetros o fenómenos no sustentados.

\subsection{Estado actual de la reconstrucción}

La implementación actual permite ejecutar una planta PMSM concentrada con parametrización M4 y conectarla a una cadena cerrada FOC.

El modelo actual no debe describirse como:

\begin{itemize}
    \item un modelo FEM;
    \item un modelo no lineal completo;
    \item un conjunto de mapas electromagnéticos validados;
    \item un modelo térmico completo;
    \item un Digital Twin terminado.
\end{itemize}

Su función actual es proporcionar una semilla reproducible y trazable sobre la cual puedan construirse esos niveles posteriores.

\subsection{Valor metodológico de la reconstrucción}

La reconstrucción Python permite transformar una recuperación de archivos en una metodología reproducible.

El conocimiento queda distribuido entre:

\begin{equation}
\text{documentación}
+
\text{ecuaciones}
+
\text{código}
+
\text{experimentos}
+
\text{resultados}.
\end{equation}

Esta estructura reduce la dependencia de un único archivo o herramienta.

También permite que una modificación futura pueda localizarse en términos de:

\begin{equation}
\text{fuente}
\rightarrow
\text{parámetro}
\rightarrow
\text{ecuación}
\rightarrow
\text{código}
\rightarrow
\text{resultado}.
\end{equation}

Esta trazabilidad constituye uno de los fundamentos del Digital Twin desarrollado en este trabajo.

\subsection{Cierre del capítulo}

La reconstrucción reproducible en Python convierte la formulación de la planta PMSM en una implementación explícita y versionable.

La estructura actual separa los parámetros M4, las ecuaciones de la planta, las interfaces de entrada y salida y las capas de control. Esta separación permite utilizar la misma planta como plataforma común para FOC y para las futuras estrategias MPC, ANFIS y TinyML.

La implementación M4 constituye actualmente una representación concentrada y reproducible. Los mapas electromagnéticos, el modelo no lineal, FEM y ROM representan etapas futuras de aumento de fidelidad y no se consideran capacidades ya implementadas.

La relación con MATLAB/Simscape permanece complementaria: MATLAB/Simscape aporta la referencia de implementación y evidencia histórica, mientras que Python proporciona la reconstrucción explícita, reproducible y versionable.

Con esta base queda establecida una semilla computacional sobre la cual puede desarrollarse la validación de la planta y, posteriormente, su evolución hacia el Digital Twin PMSM.

El siguiente capítulo documentará la validación de la planta, estableciendo qué propiedades pueden verificarse actualmente, qué métricas deben utilizarse y cuáles son todavía los límites de evidencia del modelo M4.
